\section{Annahmen: Phasen/Artefakte des DASC-PM}

- Einstieg über Quality Gates aus der Softwareentwicklung
- Nutzen zur Herleitung von Risk-Gates im DASC-PM

\notiz{
    Wie einleitend genannt, soll sich diese Arbeit mit der Strategieentwicklung für das Riskmanagement in Data-Science Projekten 
    entlang des DASC-PM als Grundlage befassen.
    Dafür soll das Werkzeug der, in der Softwareentwicklung üblichen, Quality-Gates als Ergänzungsvorschlag für das Prozessmodell 
    DASC eingeführt werden.
}

\notiz{
    TODO: Herleiten wie wir darauf kamen, dass diese Risk-Gates jetzt wirklich notwendig sind.
}

\notiz{TODO: Zunächst zu der Erwartungshaltung von Risk-Gates anhand von Quality-Gates in der Softwareentwicklung. 
    Quality-Gates, insbesondere solche die in CI / CD Pipelines genutzt werden, verwenden quantitativ erzeugte 
    Kennziffern die einen gewissen Schwellwert nicht überschreiten dürfen, da sonst die Entwicklung nicht in den 
    Hauptzweig des Projektes einfließen darf. Dies ist besonders in agilen Organisationsstrukturen sinnvoll, da 
    kleinere Meilensteine pro Iteration entwickelt werden.
    Da ein Großteil der Analysemethoden die zuvor für die Risiken genutzt wurden, ebenfalls quantitativ sind, bietet 
    sich das Prinzip des Quality Gates im Kontext von Data-Science Projekten ebenfalls an.
}

\notiz{ Um die Benennung noch treffender zu formulieren und auch eine Abgrenzung zu dem Softwareentwicklungsprozess 
    in den meißten Projekten zu erzeugen, werden die Schranken hier ''Risk Gate,, genannt. }


\begin{figure}[htbp]
    \centering
    \includegraphics[width=0.8\textwidth]{RM-Extension_DASC-PM.jpeg}
    \caption{RM-Extension für das DASC-PM}
    \label{fig:rm_extension}
\end{figure}